\section{Squarefree integers}

For squarefree integers the support classes have size one, so the graph
is described directly by subsets of the prime factors. Alternating tensors
provide invariant subspaces for this subset action~\cite{fultonharris} and
give the complete squarefree classification.

\begin{theorem}
\label{thm:squarefree}
For squarefree composite $n$, the graph $G_n$ is Laplacian integral if
and only if $n$ has exactly two prime factors.
\end{theorem}
\begin{proof}
Let $t$ be the number of prime factors. For $t=2$ the graph consists of
the two ideals $(p)$ and $(q)$, with no edge, so $L=0$. Assume $t\geq3$.
The vertices are the nonempty proper subsets of $[t]$, adjacent when they
intersect. On the tensor space encoding all subsets put
\[
E=\begin{pmatrix}2&0\\0&1\end{pmatrix},\quad
F=\begin{pmatrix}1&1\\1&0\end{pmatrix},\quad
J=\begin{pmatrix}1&1\\1&1\end{pmatrix},\quad c=2^t-1.
\]
The principal submatrix of
$\Phi=cI-E^{\otimes t}+F^{\otimes t}-J^{\otimes t}$
on the graph vertices is $L(G_n)$: the diagonal entry at $S$ is
$2^t-2-2^{t-|S|}$, and the off-diagonal entry is $-1$ for intersecting
subsets and zero otherwise.

With $e_0,e_1$ denoting the absence and presence basis, set
$w=e_0\otimes e_1-e_1\otimes e_0$.
For every $2\times2$ matrix $T$, one has
$T^{\otimes2}w=(\det T)w$. A paired factor $w$ thus contributes
$2,-1,0$ under $E,F,J$, respectively. It vanishes at the empty and
full subsets. Each subspace used below is invariant for $\Phi$ and is
supported on graph vertices, so its restriction is also a restriction of $L$.

If $t=2r+1$, $r\geq1$, use
$w^{\otimes r}\otimes\operatorname{span}\{e_0,e_1\}$.
Putting $a=2^r$ and $s=(-1)^r$, the restricted matrix is
\[
\begin{pmatrix}c-2a+s&s\\s&c-a\end{pmatrix},
\]
with discriminant $(a-s)^2+4$. The integer $h=a-s$ is at least $3$,
so $h^2<h^2+4<(h+1)^2$. Both eigenvalues are irrational.

If $t=2r$, $r\geq2$, use $r-1$ paired factors and the symmetric subspace
on the last two coordinates, with basis
$e_0e_0,(e_0e_1+e_1e_0)/\sqrt2,e_1e_1$.
Put $a=2^{r-1}$ and $s=(-1)^{r-1}$. The restriction is $cI-K$, where
\[
K=\begin{pmatrix}
4a-s&-s\sqrt2&-s\\-s\sqrt2&2a-s&0\\-s&0&a
\end{pmatrix}.
\]
The leading principal minors of $K-(a-1)I$ are
\[
3a+1-s,\quad(3a+1-s)(a+1-s)-2,\quad(3a-s)(a+1-s)-2.
\]
They are positive for $a\geq2$ and $s=\pm1$: the latter two are at
least $10$ and $8$, respectively. Sylvester's criterion gives
$\lambda_{\min}(K)>a-1$.
The principal submatrix of $K-aI$ on coordinates $1,3$ is
$\left(\begin{smallmatrix}3a-s&-s\\-s&0\end{smallmatrix}\right)$,
with determinant $-1$. A vector supported on those coordinates has a
negative Rayleigh quotient, giving $\lambda_{\min}(K)<a$.
Hence $c-\lambda_{\min}(K)$ is an eigenvalue of $L$ strictly between
the consecutive integers $c-a$ and $c-a+1$.
\end{proof}
